% META: {"id": "1NSI-PM-EX-002", "chapitre": "1NSI-PROJET-METHODES", "type_objet": "exercice", "capacites": ["BO-PREAMBULE-DEMARCHE-DE-PROJET"], "parcours": 1, "competences": ["analyser", "modeliser"], "duree_min": 10, "mode_creation": "ex_nihilo", "sources_inspiration": [], "fichier_tex": "chapitres/1NSI-PROJET-METHODES/exercices/1NSI-PM-EX-002.tex", "corrige_tex": "chapitres/1NSI-PROJET-METHODES/corriges/1NSI-PM-CO-002.tex", "status": "verified"}
% BEGIN-VERIFY
% jalons = [
%     {"nom": "Specification du jeu", "termine": True},
%     {"nom": "Deplacement du personnage", "termine": True},
%     {"nom": "Gestion des collisions", "termine": True},
%     {"nom": "Score et niveaux", "termine": False},
%     {"nom": "Tests et equilibrage", "termine": False},
%     {"nom": "Presentation finale", "termine": False},
% ]
% def avancement(jalons):
%     return sum(1 for j in jalons if j["termine"]) / len(jalons) * 100
% def prochain_jalon(jalons):
%     for j in jalons:
%         if not j["termine"]:
%             return j["nom"]
%     return None
% assert avancement(jalons) == 50.0
% assert prochain_jalon(jalons) == "Score et niveaux"
% END-VERIFY
\begin{exercice}{1NSI-PM-EX-002}{1}{10}
Une équipe de $3$ élèves développe un petit jeu vidéo, découpé en $6$ jalons :
\begin{python}
jalons = [
    {"nom": "Specification du jeu", "termine": True},
    {"nom": "Deplacement du personnage", "termine": True},
    {"nom": "Gestion des collisions", "termine": True},
    {"nom": "Score et niveaux", "termine": False},
    {"nom": "Tests et equilibrage", "termine": False},
    {"nom": "Presentation finale", "termine": False},
]
\end{python}
\begin{enumerate}
  \item En utilisant \lstinline{avancement} du cours, calculer le pourcentage d'avancement du projet.
  \item En utilisant \lstinline{prochain_jalon} du cours, déterminer sur quel jalon l'équipe doit désormais se concentrer.
  \item L'équipe est à mi-parcours de l'horaire prévu pour ce projet, et son avancement est justement de $50\,\%$. Les $3$ jalons restants (score/niveaux, tests, présentation) te semblent-ils du même volume de travail que les $3$ jalons déjà réalisés ? Pourquoi ce simple pourcentage peut-il être trompeur ?
\end{enumerate}
\end{exercice}
